\section*{अध्याय \ref{ch:b2:lab-techniques-2} --- प्रयोगशाला तकनीकें II: व्यवहार में संश्लेषण और विश्लेषण}

\begin{solution}{exo:b2:lab-techniques-2:1}
बर्फ़ और जल; बर्फ़ और नमक; शुष्क बर्फ़ और ऐसीटोन (जैसे \cref{ch:b2:enolates-aldol} के ईनोलेटों के लिए)।
\end{solution}

\begin{solution}{exo:b2:lab-techniques-2:2}
निचली परत। जल की कुछ बूँदें डालिए: वे जलीय परत से मिल जाती हैं, जो यहाँ ऊपरी है।
\end{solution}

\begin{solution}{exo:b2:lab-techniques-2:3}
पहले भाग जल ग्रहण करते हुए गुच्छे बनाते हैं; जब नया मिलाया गया ठोस ढीला रहे और स्वतंत्र रूप से घूमे,
तब विलयन शुष्क है।
\end{solution}

\begin{solution}{exo:b2:lab-techniques-2:4}
स्तंभ के लिए $R_f$ बहुत अधिक है (लगभग 0.3 चाहिए) और दोनों धब्बे निकट हैं: कम ध्रुवी
निक्षालक प्रयुक्त कीजिए, अधिक हेक्सेन (उदाहरण के लिए 9\,:\,1), जो दोनों $R_f$ को घटाता है और स्तंभ आयतनों में
उनका पृथक्करण बढ़ाता है।
\end{solution}

\begin{solution}{exo:b2:lab-techniques-2:5}
\begin{itemize}
  \item सोडियम हाइड्रोजनकार्बोनेट (pH 8.3 $>$ 4.2 + 2): जल में बेन्ज़ोएट; अम्लीय बनाइए, बेन्ज़ोइक अम्ल निष्कर्षित कीजिए।
  \item सोडियम हाइड्रॉक्साइड (pH $>$ 12 $>$ 9.99 + 2): जल में फ़ीनॉक्साइड; अम्लीय बनाइए, फ़ीनॉल निष्कर्षित कीजिए।
  \item हाइड्रोक्लोरिक अम्ल (pH लगभग 1 $<$ 4.6 $-$ 2): जल में ऐनिलीनियम; क्षारीय बनाइए, ऐनिलीन निष्कर्षित कीजिए।
  \item नैफ़्थलीन ईथर में रहता है: सुखाइए, वाष्पित कीजिए।
\end{itemize}
हाइड्रोजनकार्बोनेट धुलाई पहले आनी चाहिए: हाइड्रॉक्साइड अम्ल और फ़ीनॉल का एक साथ निष्कर्षण कर लेता।
% ledger: ph:NaHCO3, pka:benzoic, pka:phenol, pka:anilinium
\end{solution}

\begin{solution}{exo:b2:lab-techniques-2:6}
$1/T = 1/T_b - (R/\Delta_{\mathrm{vap}}H)\ln(p/p^\circ)$: ब्रोमोबेन्ज़ीन लगभग \qty{49}{\degreeCelsius},
बेन्ज़ैल्डिहाइड लगभग \qty{68}{\degreeCelsius} (\cref{prop:b2:lab-techniques-2:reduced-pressure} का चित्र)।
\end{solution}

\begin{solution}{exo:b2:lab-techniques-2:7}
$0.050 \times 200\,000/400 = \qty{25}{K}$।
\end{solution}

\begin{solution}{exo:b2:lab-techniques-2:8}
हर यौगिक की अपनी अनुक्रिया होती है: क्षेत्रफल प्रतिशत द्रव्यमान प्रतिशत तभी है जब \omterm{def:b2:chromatography:quantitative}{अनुक्रिया गुणांक}
समान हों। NMR: $0.12/3 = 0.040$~mol अशुद्धि प्रति mol उत्पाद, अर्थात् उत्पाद के \qty{200}{g} प्रति
$0.040 \times 100 = \qty{4}{g}$; $w = 200/204 \approx 98.0~\%$।
\end{solution}

\begin{solution}{exo:b2:lab-techniques-2:9}
सैद्धांतिक द्रव्यमान $0.0300 \times 164 = \qty{4.92}{g}$। असंशोधित $4.10/4.92 \approx 83.3~\%$; संशोधित
$0.950 \times 4.10/4.92 \approx 79.2~\%$।
\end{solution}

\begin{solution}{exo:b2:lab-techniques-2:10}
बेन्ज़ीन का अर्थ है कि अभिकर्मक बना और फिर प्रोटॉनित हो गया: काँच-उपकरणों, विलायक या
ऐल्डिहाइड में जल; किसी रिसाव से प्रवेश करती वायु (नमी); आंशिक रूप से बेन्ज़ोइक अम्ल में ऑक्सीकृत ऐल्डिहाइड, जिसका
अम्लीय हाइड्रोजन अभिकर्मक को नष्ट कर देता है। उपाय: ओवन में सुखाए काँच-उपकरण, शुष्क विलायक, ताज़ा आसवित
ऐल्डिहाइड, बुदबुदक पर जाँचा गया नाइट्रोजन अधिदाब। यदि अभिकर्मक लगभग बना ही न हो तो यह
आरंभन की विफलता (ऑक्साइड से ढकी मैग्नीशियम सतह) की ओर संकेत करता, जिससे ताज़ी या सक्रियित छीलन द्वारा बचा जाता है।
\end{solution}

\begin{solution}{exo:b2:lab-techniques-2:11}
\qty{150}{W} पर \qty{270}{kJ}: कम-से-कम $270\,000/150 = \qty{1800}{s}$, आधा घंटा, और केवल तब जब
हैलाइड मिलाते ही उतनी तेज़ी से अभिक्रिया करे। इससे तेज़ होने पर ऊष्मा मुक्ति शीतन से अधिक हो जाती है: ताप बढ़ता है,
या अभिकर्मक संचित होता है और बाद में एक साथ अभिक्रिया कर सकता है (\cref{prop:b2:lab-techniques-2:accumulation})।
\end{solution}

\begin{solution}{exo:b2:lab-techniques-2:12}
पुनःक्रिस्टलन: 99~\% पर $13.0 \times 0.80 = \qty{10.4}{g}$। स्तंभ: 99.5~\% पर $13.0 \times 0.95 =
\qty{12.4}{g}$, परंतु \qty{1.5}{L} विलायक जिसे ख़रीदना, वाष्पित करना और निपटाना पड़ता है। स्तंभ लगभग
\qty{2}{g} अधिक और थोड़ा अधिक शुद्ध उत्पाद देता है; पुनःक्रिस्टलन सरल, सस्ता और कम
अपव्ययी है, और जब उसकी शुद्धता पर्याप्त हो तो उसे वरीयता दी जाती है।
\end{solution}

\begin{solution}{pb:b2:lab-techniques-2:1}
\textbf{1.} एथॉक्सीएथेन: अत्यंत ज्वलनशील, हानिकारक, परॉक्साइड बनाता है। ब्रोमोबेन्ज़ीन: ज्वलनशील, विषैला,
जलीय जीवों के लिए विषैला। मैग्नीशियम: ज्वलनशील ठोस, जल के साथ हाइड्रोजन मुक्त करता है।
\textbf{2.} \ce{C6H5MgBr + H2O -> C6H6 + Mg(OH)Br}: जल की हर बूँद अभिकर्मक को नष्ट करती है; ओवन में सुखाना
काँच पर जल की परत को हटा देता है।
\textbf{3.} नमी और ऑक्सीजन को बाहर रखने के लिए, जो भी अभिकर्मक को नष्ट करती है।
\textbf{4.} नाइट्रोजन का प्रवाह; यह निकास से वायु के प्रवेश को रोकता है।
\textbf{5.} यह अभिकर्मक से अभिक्रिया नहीं करता, इसके ऑक्सीजन के एकाकी युग्म मैग्नीशियम को स्थायी करते हैं, और इसका
निम्न क्वथनांक (\qty{307.7}{K}) पश्चवाह को ताप सीमित करने देता है।
\textbf{6.} $15.7/157.01 = \qty{0.1000}{mol}$ ब्रोमोबेन्ज़ीन; $2.67/24.305 = \qty{0.1099}{mol}$ मैग्नीशियम,
10~\% आधिक्य में।
\textbf{7.} $0.1000 \times 270 = \qty{27.0}{kJ}$।
\textbf{8.} ताकि ऊष्मा उसी दर से मुक्त हो जिस दर से पश्चवाह और कुंड उसे हटाते हैं।
\textbf{9.} $0.0200 \times 270 = \qty{5.40}{kJ}$।
\textbf{10.} $5400/170 \approx \qty{32}{K}$।
\textbf{11.} लगभग \qty{52}{\degreeCelsius}, एथॉक्सीएथेन के क्वथनांक (\qty{34.6}{\degreeCelsius}) से बहुत ऊपर:
प्रचंड क्वथन जो संघनित्र पर हावी होकर ज्वलनशील वाष्प को
फ़्लास्क से बाहर फेंक सकता है।
\textbf{12.} ब्रोमोबेन्ज़ीन का लगभग दसवाँ भाग मिलाइए, अभिक्रिया आरंभ होने की प्रतीक्षा कीजिए (धुँधलापन, हल्का
पश्चवाह), फिर शेष को हल्के पश्चवाह की दर से मिलाइए, बर्फ़ कुंड तैयार रखते हुए।
\textbf{13.} \ce{C6H5MgBr + C6H5CHO} ऐल्कॉक्साइड \ce{(C6H5)2CHOMgBr} देता है; बेन्ज़ैल्डिहाइड,
$9.55/106.12 = \qty{0.0900}{mol}$, सीमांत है।
\textbf{14.} ऐल्कोहॉल दोहरा बेन्ज़िलिक है: प्रबल अम्ल उसे एक स्थायीकृत धनायन में बदल देता,
जिससे ईथर या प्रतिस्थापन उत्पाद बनते; अमोनियम क्लोराइड ऐल्कॉक्साइड को प्रोटॉनित करने
और मैग्नीशियम लवणों को घोलने के लिए पर्याप्त अम्लीय है।
\textbf{15.} ऊपरी परत: एथॉक्सीएथेन जल से कम सघन है।
\textbf{16.} यह ईथर में घुले जल का अधिकांश भाग हटा देता है।
\textbf{17.} मुक्त-प्रवाही होने तक मैग्नीशियम सल्फ़ेट, निस्यंदन, फिर घटे हुए दाब पर घूर्णी
वाष्पित्र।
\textbf{18.} फ़ेनिलमैग्नीशियम ब्रोमाइड अनभिक्रियित ब्रोमोबेन्ज़ीन से युग्मित होता है; ब्रोमोबेन्ज़ीन की उच्च स्थानीय सांद्रता
और उच्च ताप इसे बढ़ावा देते हैं, जो धीमे योजन का एक और कारण है।
\textbf{19.} बाइफ़ेनिल ($R_f$ 0.85), कम ध्रुवी।
\textbf{20.} बाइफ़ेनिल के लिए लगभग $1/0.85 \approx 1.2$ स्तंभ आयतन, ऐल्कोहॉल के लिए $1/0.25 = 4$।
\textbf{21.} ऐल्कोहॉल प्रति CH 10 \omterm{def:b2:huckel:aromatic}{ऐरोमैटिक} प्रोटॉन देता है; अतिरिक्त $10.90 - 10.00 = 0.90$
बाइफ़ेनिल से आता है, प्रति अणु 10 प्रोटॉन: प्रति mol ऐल्कोहॉल 0.090 mol बाइफ़ेनिल।
\textbf{22.} $w = 184.24/(184.24 + 0.090 \times 154.21) \approx 0.930$।
\textbf{23.} $12.50 \times 0.930 \approx \qty{11.62}{g}$।
\textbf{24.} $0.0900 \times 184.24 \approx \qty{16.58}{g}$।
\textbf{25.} $12.50/16.58 \approx 75.4~\%$।
\textbf{26.} शुद्धता-संशोधित, $11.62/16.58 \approx \boldsymbol{70~\%}$ (70.1~\%)।
\end{solution}
